% group5_qed_gauge_synthesis.tex
% GeoVac Group 5 Synthesis: QED / Gauge / Gravity Arc
% Status: Project archive synthesis, June 2026
% Audience: HEP / gauge-theory / NCG-adjacent readers
% Spine: forced gauge structure (25,30,41) -> graph-native QED (33,28)
%        -> the alpha observation (2) -> bound-state precision QED (36)
%        -> spectral-action gravity sketch (51) -> honest scope

\documentclass[aps,pra,twocolumn,superscriptaddress,longbibliography]{revtex4-2}

\usepackage{amsmath,amssymb,amsthm,graphicx,xcolor,bm,braket,booktabs}
\usepackage{hyperref}

\newtheorem{theorem}{Theorem}
\newtheorem{proposition}{Proposition}
\newtheorem{corollary}{Corollary}
\newtheorem{observation}{Observation}

\begin{document}

\title{Gauge Structure, QED, and Gravity from the Spectral Geometry of
$S^3$:\\ A Synthesis of the GeoVac QED/Gauge Arc}
\thanks{Group~5 synthesis in the Geometric Vacuum series}

\author{J.~Loutey}
\affiliation{Independent Researcher, Kent, Washington}

\date{June 30, 2026}

%% ========================================================================
%% ABSTRACT
%% ========================================================================

\begin{abstract}
The Geometric Vacuum (GeoVac) framework discretizes the hydrogen
bound-state manifold as the unit three-sphere $S^3$ via Fock's conformal
projection.  This synthesis assembles the eight Group~5
papers---gauge structure, graph-native QED, precision bound-state QED,
and a spectral-action gravity sketch---into one narrative around a single
organizing distinction:\ \emph{the gauge group structure is largely
forced by the spectral geometry of $S^3$ and its Hopf fibration
$S^1\!\to\!S^3\!\to\!S^2$, while the couplings are
calibration/observation, not derived}.  The graph carries a $U(1)$
Wilson/Hodge lattice gauge structure~\cite{loutey_paper25}, a non-abelian
$\mathrm{SU}(2)$ Wilson theory because $S^3=\mathrm{SU}(2)$ as a
manifold~\cite{loutey_paper30}, and---on the cross-$\ell$ Dirac
graph---a $U(1)$ structurally compatible with three-dimensional compact
electrodynamics~\cite{loutey_paper41}; the natural gauge reach is
$U(1)\times\mathrm{SU}(2)\times\mathrm{SU}(3)$, with $\mathrm{SU}(3)$
entering only as a pure-gauge theory whose natural matter coupling is
obstructed.  The framework's QED is genuinely graph-native in its
selection-rule and period content:\ the eight perturbative selection
rules partition $1{+}6{+}1$ by which graph ingredient recovers
them~\cite{loutey_paper33}, and the $S^3$ Dirac spectral sums reproduce
the flat-space QED transcendentals ($\zeta(3)$, $\zeta(5)$, Catalan~$G$,
$\zeta(5,3)$) through a depth-filtered period
tower~\cite{loutey_paper28}.  The fine-structure constant satisfies the
cubic $\alpha^3-K\alpha+1=0$ with $K=\pi(B+F-\Delta)$ assembled from
three spectral invariants of the Hopf bundle, each with an independently
established spectral home; this
is recorded strictly as an \textbf{Observation} (agreement
$8.8\times10^{-8}$, twelve single-mechanism derivations eliminated),
never as a derivation~\cite{loutey_paper2}.  Bound-state QED reaches
sub-percent one-loop Lamb-shift closure with no GeoVac-specific
fit~\cite{loutey_paper36}, and a spectral-action gravity sketch yields an
\emph{exactly} two-term Einstein--Hilbert action together with a set of
honest negatives that bound a still-open discrete-substrate
program~\cite{loutey_paper51}.  The discipline throughout matches the
rest of the corpus:\ every load-bearing claim wears its tier, the
structural results are foregrounded, and the couplings are positioned as
the free, calibration side of the construction.
\end{abstract}

\maketitle

%% ========================================================================
\section{Introduction: forced structure, free couplings}
%% ========================================================================

The GeoVac program establishes that the discrete graph Laplacian for
hydrogen is mathematically equivalent to the Laplace--Beltrami operator
on the unit three-sphere~$S^3$, related to the Schr\"odinger equation by
Fock's 1935 conformal projection~\cite{loutey_paper7}.  The Group~5 arc
asks what the \emph{same} spectral geometry says about the rest of
electromagnetism, the gauge sector, and---most speculatively---gravity.
A single distinction organizes the answer, and it is the through-line of
this synthesis.

\emph{The gauge group structure is largely forced.}  The three-sphere
carries a canonical principal-bundle structure, the Hopf fibration
$S^1\!\to\!S^3\!\to\!S^2$, and $S^3$ \emph{is} the group manifold
$\mathrm{SU}(2)$.  These two facts are not GeoVac inputs; they are
geometry.  Reading the GeoVac graph through its incidence matrix turns
them into an abelian and then a non-abelian Wilson lattice gauge theory
on the same nodes and edges that carry the Coulomb spectrum
(\S\ref{sec:gauge}).  The natural gauge reach saturates at
$U(1)\times\mathrm{SU}(2)\times\mathrm{SU}(3)$, and where it stops is
itself structural:\ $\mathrm{SU}(3)$ admits only a pure-gauge Wilson
theory, its natural matter coupling obstructed at the Clebsch--Gordan
intertwiner level (a sharpening of the Coulomb/harmonic-oscillator
asymmetry of Paper~24~\cite{loutey_paper24}).

\emph{The couplings are not.}  The fine-structure constant, the QED
multi-loop constants, and the cutoff function of the gravity sector are
all on the other side of the ledger:\ calibration or observation, in the
two-layer language of Paper~34~\cite{loutey_paper34}.  The framework's
most precise number---$\alpha$ reproduced to $8.8\times10^{-8}$---is held
explicitly at \textbf{Observation} tier, with no claim of derivation
(\S\ref{sec:alpha}).  This is the same discipline that runs through the
quantum-chemistry and quantum-computing arcs:\ lead with the structural
result, name the calibration boundary, and let the tier of each claim be
visible in the prose.

Following the framework's positioning convention, the discrete graph
topology and the continuous Hopf/spectral-action geometry are treated as
\emph{dual descriptions} connected by proven maps;\ no ontological
priority is asserted for either, and the synthesis leads with concrete
structural content rather than interpretive claims about the nature of
electromagnetism or gravity.

The eight papers form a loose dependency order.  Papers~25, 30, and~41
build the gauge structure.  Papers~33 and~28 develop graph-native QED:\
its selection rules and its perturbative period content.  Paper~2 reports
the $\alpha$ observation.  Paper~36 turns the Dirac-$S^3$ machinery onto
precision bound-state observables.  Paper~51 sketches the spectral-action
gravity sector and---characteristically---records what did not work.

%% ========================================================================
\section{Gauge structure from the Hopf bundle}\label{sec:gauge}
%% ========================================================================

\paragraph{The Hodge/Wilson dictionary on the graph (Paper~25).}
With $B$ the node--edge incidence matrix of the GeoVac graph, the node
Laplacian $L_0=BB^{\top}$ is the matter propagator (it converges to
$-\Delta_{S^3}$~\cite{loutey_paper7}) and the edge Laplacian
$L_1=B^{\top}B$ is the discrete Hodge-1 Laplacian whose kernel counts
cycles---the photon propagator in temporal gauge.  Assigning a $U(1)$
connection $U_e=e^{i\theta_e}$ to each edge and collapsing the magnetic
label $m_\ell$ over its $2\ell+1$ Hopf fiber reproduces the $S^2$ base
graph exactly.  None of the underlying mathematics is new (graph Hodge
theory, Wilson lattice gauge theory, the Marcolli--van~Suijlekom
gauge-network spectral triple~\cite{marcolli_vs}); what Paper~25 claims is
the \emph{synthesis}---one physically motivated discretization on which
matter, an abelian connection, and the Hopf base coexist.  The paper is
explicit about its tier:\ ``an observational synthesis paper. No new
experimental or computational prediction is derived,'' and the
edge-Laplacian spectrum of a finite Hopf graph is necessarily
algebraic ($\pi$-free), so no transcendental can hide in the gauge
1-forms themselves.

\paragraph{$\mathrm{SU}(2)$ because $S^3=\mathrm{SU}(2)$ (Paper~30).}
The rank-1 non-abelian case is special:\ the manifold \emph{is} the gauge
group.  With links $U_e\in\mathrm{SU}(2)$, plaquettes the primitive
non-backtracking (Ihara) cycles, and Wilson action
$S_W=\beta\sum_P\bigl(1-\tfrac12\,\mathrm{Re}\,\mathrm{tr}\,U_P\bigr)$,
the construction is a gauge-invariant, Haar-normalizable finite theory.
Two results carry real weight.  The maximal-torus (Cartan) reduction
reproduces Paper~25's $U(1)$ action to machine precision
($|S_W-S_{U(1)}|<10^{-12}$)---stated as a \textbf{proposition with a
proof}---so Paper~25 is exactly the abelian sub-sector of a genuinely
non-abelian theory.  The second-order weak-coupling expansion produces
the co-exact (curl) plaquette form $B_2B_2^{\top}$ --- \emph{not} a
multiple of $L_1$; it is the missing half that completes Paper~25's
$L_1=B^{\top}B$ to the full discrete Hodge-1 Laplacian (Paper~30
Prop.~2, corrected 2026-07-03) --- with non-abelian couplings first
entering at order $A^4$.  The kinetic normalization is
universal,
\begin{equation}
\text{coeff} = \frac{1}{4N_c}
\;\Longrightarrow\;
\tfrac18\ (\mathrm{SU}(2)),\quad
\tfrac1{12}\ (\mathrm{SU}(3)),
\end{equation}
verified symbolically and numerically.  The scope is explicitly bounded:\
this is a compact-graph construction, \emph{not} a Yang--Mills mass-gap
or confinement claim on $\mathbb{R}^4$; the Monte-Carlo plaquette data is
leading-order and agrees with the character expansion only to the
$\sim\!20\%$ level at intermediate $\beta$.

\paragraph{Where the gauge reach stops.}
The non-abelian story does not continue to $\mathrm{SU}(3)$ as a
\emph{natural} matter-coupled theory.  On the $S^5$ Bargmann--Segal
lattice (the harmonic-oscillator analog of $S^3$) a pure-gauge Wilson
$\mathrm{SU}(3)$ is well defined and machine-precision verified (the same
$1/(4N_c)=1/12$ coefficient, Cartan reduction over $7{,}765$ plaquettes
to $<1.5\times10^{-11}$), but the natural matter sector cannot be coupled
in the strict Wilson sense:\ the Clebsch--Gordan intertwiner obstruction
re-appears at the matter level, and a least-squares fit of the
$m_\ell$-quotient to $\mathbb{CP}^2$ fails uniformly ($40.8\%$ least-squares
residual, with a convention-robust $\geq 33\%$ floor; Paper~25 \S VII).  This is the gauge-side face of the Coulomb/HO asymmetry
catalogued in Paper~24~\cite{loutey_paper24}:\ $U(1)$ transfers
universally, $\mathrm{SU}(2)$ is forced by $S^3=\mathrm{SU}(2)$, and
$\mathrm{SU}(3)$ is gauge-natural but matter-obstructed.  The cumulative
reading is that the natural gauge content of the construction saturates
at $U(1)\times\mathrm{SU}(2)\times\mathrm{SU}(3)$---the structural reach,
not a derivation of the Standard-Model gauge group, and consistent with
the spectral-triple gauge analysis of Paper~32~\cite{loutey_paper32}.

\paragraph{The Dirac Rule-B graph and 3D compact $U(1)$ (Paper~41).}
Paper~25's scalar Wilson $U(1)$ has a structural limitation:\ on the
scalar Fock graph the gauge theory factorizes into disjoint per-$\ell$
two-dimensional grids, so it is really 2D compact $U(1)$ with no genuine
three-dimensional phase structure.  The cross-$\ell$ Dirac ``Rule~B''
graph (the $E1$ dipole adjacency, $|\Delta\ell|=1$) repairs this:\ every
edge flips $\ell$-parity, the graph is a single connected component, and
its spectral dimension climbs $1.86\to2.27\to2.54$ toward~3 as
$n_{\max}=3\to5$.  Wilson $U(1)$ on Rule~B is then tested against seven
independent structural witnesses (spectral dimension, Migdal--Kadanoff RG
flow, Gaussian perimeter law, strong-coupling area law, a DeGrand--Toussaint
monopole-density signature decaying over five decades, full
Monte-Carlo Wilson loops, and finite-$T$ Polyakov loops on
$G_B\times C_{N_t}$).  Six pass cleanly and one (full-MC Wilson loops) is
a ``nuanced weak pass.''  The verdict is stated at exactly the tier the
evidence supports:\ the theory is \emph{structurally compatible} with 3D
compact $U(1)$ (the permanent-confinement, no-finite-$T$-deconfinement
universality class) and inconsistent with 4D---\textbf{explicitly not}
``uniquely identified with 3D up to universality,'' and with the
$n_{\max}\!\to\!\infty$ continuum limit left uncomputed.

%% ========================================================================
\section{Graph-native QED: selection rules and periods}\label{sec:qed}
%% ========================================================================

\begin{table}[t]
\centering
\begin{tabular}{lcl}
\hline\hline
Construction & Count & Tier reading \\
\hline
Scalar graph $+$ scalar photon      & $1/8$ & I: graph-native (Gaunt) \\
Native Dirac graph (scalar photon)  & $4/8$ & $+$ spinor labels \\
Vector photon $+$ scalar electrons  & $7/8$ & II: $+6$ at $1/(4\pi)$/loop \\
Vector photon $+$ Dirac electrons   & $8/8$ & III: $+$ Furry, free \\
\hline\hline
\end{tabular}
\caption{The $1{+}6{+}1$ partition of the eight QED selection rules by
recoverability on the GeoVac construction.  The word ``scalar'' is
overloaded:\ the $1/8$ row is the scalar \emph{photon} on the scalar
graph;\ the $7/8$ row is scalar \emph{electrons} with the \emph{vector}
photon.}
\label{tab:partition}
\end{table}

The QED arc asks how much of perturbative quantum electrodynamics is
visible on the bare combinatorial graph, before any continuum projection.
Two complementary answers emerge:\ a clean partition of the selection
rules, and a clean accounting of the perturbative transcendentals.

\paragraph{The $1{+}6{+}1$ partition of selection rules (Paper~33).}
The eight standard perturbative-QED selection rules partition into three
tiers by \emph{which} ingredient of the construction is needed to recover
each (Table~\ref{tab:partition}).

On the scalar Fock graph with a scalar ($1$-cochain) photon, exactly
\emph{one} rule survives---the Gaunt/Clebsch--Gordan angular-momentum
sparsity baked into the discrete edge set (graph-native, Tier~I, free).
Promoting the photon to a vector harmonic carrying explicit angular
labels $(q,m_q)$ with Wigner-$3j$ vertex coupling recovers
\emph{six} angular-momentum rules at once, at a cost of exactly
$1/(4\pi)$ per loop---identified as the $S^2$ Weyl exchange constant of
the Hopf base, an entry in Paper~18's calibration tier~\cite{loutey_paper18}
(Tier~II).  The eighth rule, charge conjugation / Furry, is recovered for
\emph{Dirac} electrons at zero cost by a single-particle spinor phase
constraint $V(a,a,q,m_q)=0$ (proven, sympy-exact for six cases), a
kinematic identity distinct from second-quantized $C$-symmetry
(Tier~III).  The load-bearing structural statement is the dividing line:\
the photon must carry angular quantum numbers $(L,M_L)$, not merely a
spin index at the vertex, to recover the six angular rules---which is
precisely why those six sit on the calibration side and only the Gaunt
rule is graph-native.

\paragraph{The perturbative period tower (Paper~28).}
The $S^3$ Dirac spectral geometry (Camporesi--Higuchi spectrum
$|\lambda_n|=n+\tfrac32$, Hodge-1 photon, $\mathrm{SO}(4)$ vertex rules)
produces the \emph{same} transcendentals as flat-space perturbative QED
through a purely spectral mechanism organized by operator order and
vertex topology.  Three results anchor the arc, at theorem/proposition tier (the realized-depth tower is a Proposition):
\begin{itemize}
\item \textbf{T9 (two-term theorem).}  The $D^2$ spectral zeta is a
two-term $\pi^{\text{even}}$ polynomial with rational coefficients at
every integer $s$---no odd-zeta content; operator order (even vs.\ odd
$s$) is the primary transcendental discriminant.
\item \textbf{$\chi_{-4}$ bridge (Theorem~3).}  A spectral-zeta-to-Dirichlet-$L$
identity, $D_{\text{even}}(s)-D_{\text{odd}}(s)=2^{\,s-1}\bigl(\beta(s)-\beta(s-2)\bigr)$
with $\beta(s)=L(s,\chi_{-4})$, in which vertex parity acts as the
Dirichlet character; proven, and verified symbolically and by PSLQ to
100 digits.  It is consistent with but does not enforce the Riemann
hypothesis for $\beta$.
\item \textbf{Realized-depth tower.}  Loop order bounds the multiple-zeta
depth but does not automatically mint new irreducibles; the realized
depth grows strictly more slowly, $\le k-1$.  The two-loop sunset
constant reduces entirely,
$S_{\min}\in\mathbb{Q}[\pi^2,\ln2,\zeta(3),\zeta(5)]$ (depth $\le1$);
the three-loop chain $S^{(3)}=31.5726\ldots$ first introduces a genuine
depth-2 period, $\zeta(5,3)$, as its \emph{only} irreducible beyond the
single-zeta ring.
\end{itemize}
On the finite graph, all of these quantities are $\pi$-free rationals or
algebraic numbers---the transcendentals enter only through the infinite
spectral sums of the continuum projection.  The continuum self-energy
zero is correspondingly \emph{broken} on the graph (the pendant-edge
proposition gives $\Sigma_{\text{graph}}(\text{GS})=2(n_{\max}-1)/n_{\max}\to2\neq0$),
and is recovered by either of two routes: the vector-photon angular
labels of Paper~33, or the plaquette (co-exact) transverse propagator
of Papers~28/30.  A spectral-sum anomalous moment
$F_2/[\alpha/2\pi]=1.084$ lands within $0.4\%$ of the Parker--Toms curved-space
value---reported as a numerical match, not a derivation.

%% ========================================================================
\section{The fine-structure-constant observation}\label{sec:alpha}
%% ========================================================================

The arc's most precise---and most carefully bounded---result is
Paper~2's.  The fine-structure constant satisfies the depressed cubic
\begin{equation}
\alpha^3-K\alpha+1=0,\qquad
K=\pi\!\left(B+F-\Delta\right)
 =\pi\!\left(42+\frac{\pi^2}{6}-\frac{1}{40}\right),
\label{eq:K}
\end{equation}
where the three coefficients are spectral invariants of the Hopf bundle
components:\ $B=42$ is the degeneracy-weighted Casimir trace of the $S^2$
base (with the selection principle $B/N=\dim S^3=3$ proven algebraically
to fix the cutoff $n_{\max}=3$ uniquely), $F=\zeta(2)=\pi^2/6$ is the
spectral zeta of the $S^1$ fiber (equivalently the Fock-degeneracy
Dirichlet series at the packing exponent $d_{\max}=4$~\cite{loutey_paper0}),
and $\Delta=1/40$ is an $S^3$ boundary term equal to the reciprocal of
the single-level Dirac degeneracy $g_3^{\text{Dirac}}=2(n{+}1)(n{+}2)|_{n=3}=40$.
The physical root gives $\alpha^{-1}=137.036011$, matching CODATA to a
relative error of $8.8\times10^{-8}$ with zero free parameters; a
combinatorial search over $1.92\times10^9$ comparable formulas returns a
$p$-value of $5.2\times10^{-9}$.

The status of Eq.~\eqref{eq:K} must be stated exactly, and it is a
\textbf{hard prohibition} of the framework that it be stated no more
strongly.  The combination rule $K=\pi(B+F-\Delta)$ is an
\textbf{Observation}:\ an empirical numerical coincidence with structural
support, \emph{never} a derivation, conjecture, prediction, or theorem.
Each ingredient has an \emph{independently derived} spectral home---a
``three-homes'' theorem establishes that $B$, $F$, $\Delta$ live in two
distinct sectors of $S^3$ (scalar and spinor), are three different object
types (finite trace, infinite Dirichlet series, single-level degeneracy),
and that no spectral construction on $S^3$ generates all three at
once---but their \emph{combination} is not derived.  A multi-phase
structural-decomposition program eliminated twelve single-mechanism
derivations of the rule, leaving the cross-sector sum a genuine open
question:\ in the paper's own assessment, the geometric arena (Link~1)
and the per-component invariants (Link~2) are largely established, while the
additive combination (Link~3) is ``not established.''  The synthesis
inherits this verdict verbatim:\ what is structural here is each spectral
home; what is observational is that they sum to $\alpha^{-1}/\pi$.

%% ========================================================================
\section{Bound-state precision QED}\label{sec:bsqed}
%% ========================================================================

Paper~36 turns the Dirac-$S^3$ infrastructure onto Lamb-shift-class
observables, where the answer is known independently and the test is one
of internal consistency rather than prediction.  The hydrogen
$2S_{1/2}$--$2P_{1/2}$ Lamb shift comes out at $1052.19$~MHz against the
experimental $1057.845$~MHz---a $-0.534\%$ one-loop residual---using only
the bare graph plus three standard projections (Fock, the
Camporesi--Higuchi spinor lift, and the Drake--Swainson asymptotic
subtraction, which is Paper~34's thirteenth catalogued
projection~\cite{loutey_paper34}), with \emph{no} GeoVac-specific fitted
parameter.  Companion systems land comparably:\ muonic hydrogen at
$-0.10\%$ and the muonium $2S$--$2P$ Lamb shift at $+0.013\%$ (the
cleanest match in the set).

The honest reading is that this is \emph{observation-side} physics in the
two-layer sense.  The framework genuinely \emph{derives} the spectral-sum
structure---the Bethe logarithms, computed natively in the
acceleration form, agree with Drake--Swainson to better than $1\%$---but
\emph{imports} the QED constants (the $10/9$ one-loop self-energy
coefficient, the Uehling kernel) from the standard Eides convention; in
the Paper~34 language, Layer~1 (the bare Dirac-$S^3$ spectrum) is
$\pi$-free, and every $\pi$ enters through a Layer-2 projection.  The
two-loop sector is the explicit boundary:\ the bare iterated
spectral action reproduces the flat-space two-loop \emph{divergence}
faithfully but does not generate renormalization counterterms, so finite
multi-loop predictions remain external input.  The result is therefore
positioned as one-loop closure with the two-loop closure descoped to a
renormalization extension---a precision test passed, not a new QED
calculation.

%% ========================================================================
\section{The spectral-action gravity sketch}\label{sec:gravity}
%% ========================================================================

The most speculative paper in the arc, Paper~51, asks whether the
Fock-projected $S^3$ supports a sensible gravity reading at the
Chamseddine--Connes spectral-action level~\cite{chamseddine_connes}.  It
is positioned, correctly and throughout, as an \emph{exploratory,
structural-sketch} program---``with the standard caveats of the
Connes--Chamseddine program''---not as established physics, and it is
unusually candid about what failed.

One result is solid.  A spectral-zeta identity on the Camporesi--Higuchi
spinor spectrum, $\zeta_{\text{unit}}(-k)=0$ for all $k\ge0$ (proven via
a Bernoulli identity $B_{2k+1}(3/2)=(2k+1)/4^k$), \emph{forces} the
spectral action on $S^3$ to be \emph{exactly two-term}---Einstein--Hilbert
plus a cosmological constant, with no higher-curvature corrections at any
order.  This is a genuine theorem, and it supplies the algebraic
mechanism behind the ``remarkable cancellations'' noted by
Chamseddine--Connes; it is spinor-bundle-specific (the scalar and
transverse-traceless sectors retain the full infinite Seeley--DeWitt
series).  From it follow, at the spectral-action level, a Bekenstein--Hawking
entropy $S_{\text{BH}}=A\Lambda^2/(12\pi)$ (continuum replica
derivation), an effective Newton constant $G_{\text{eff}}=6\pi/\Lambda^2$,
a de~Sitter vacuum, a proven Fierz--Pauli ``$J$-blindness'' theorem, and
a proven (Layer-2, rate-level) replica-weight-harmless theorem.

The negatives are the point, and they are stated plainly:
\begin{itemize}
\item the conical-defect tip coefficient $A=1/(24\pi)$ was \emph{supplied
analytically}, not extracted---all three discrete-substrate
discretization schemes undershoot the target;
\item the M\"obius factor $\alpha/(2\alpha-1)$, once read as a continuum
$\alpha>1$ ``mechanism,'' is a \emph{finite-lattice-spacing artifact}:\
as the radial mesh refines, the recovery climbs back toward the plain
Sommerfeld--Cheeger value, so there is no exotic mechanism to find;
\item no cutoff function cures the infrared over-count across all
$\Lambda$ (clean negative);
\item the Newton-constant ``factor of~2'' is a Wald-forced bookkeeping
artifact, \emph{not} a scalar-vs-Dirac cone-coefficient difference (2D
Dirac and scalar share the same $(1/12)(1/\alpha-\alpha)$ cone term);
\item Bisognano--Wichmann wedge entanglement entropy is \emph{not} an
area law---it scales as $\sim2\log n_{\max}$ (area-law fit rejected at
$R^2=0.83$), so the entropy comes from the replica/spectral-action route,
not from wedge entanglement.
\end{itemize}
The overall verdict is that GeoVac fixes the \emph{structural skeleton} of
gravity at the spectral-action level---the two-term action is forced---while
the cutoff function remains external calibration data and the full
discrete-substrate re-derivation of $S_{\text{BH}}$ is an explicitly
open research program.  Represented at its true tier, this is a
proven two-term-exactness theorem wrapped in an honest, still-open
exploratory arc.

%% ========================================================================
\section{Honest scope and positioning}
%% ========================================================================

\paragraph{What is robust.}
The load-bearing results are structural and proven.  The Hodge/Wilson
dictionary ($L_0=BB^{\top}$ matter, $L_1=B^{\top}B$ photon) is exact, as
is the machine-precision $\mathrm{SU}(2)$ maximal-torus reduction that
makes Paper~25 the Cartan sector of Paper~30.  The QED selection-rule
partition $1{+}6{+}1$ and the perturbative period results---the T9
two-term theorem, the $\chi_{-4}$ Dirichlet-$L$ bridge, and the
realized-depth tower with $\zeta(5,3)$ as the first genuine depth-2
period---are theorems with symbolic/PSLQ backing.  The spectral-action
two-term exactness on $S^3$ is a proven theorem.  Each of the three
$\alpha$-ingredients $B$, $F$, $\Delta$ has an independently derived
spectral home.

\paragraph{What is conditional or observational.}
The fine-structure combination rule $K=\pi(B+F-\Delta)$ is an
\textbf{Observation}, not a derivation---the single most important tier
statement in the arc.  The bound-state QED matches are observation-side:\
the spectral-sum structure is derived but the QED constants are imported,
and every $\pi$ enters through a Layer-2 projection.  The Rule-B gauge
theory is ``structurally compatible with 3D compact $U(1)$,'' not
``identified'' with it, and its spectral-dimension-to-3 evidence rests on
a finite-size extrapolation.  The entire gravity sector is a
spectral-action-level structural sketch with an open discrete-substrate
program.

\paragraph{What was tried and did not work.}
The arc is bounded by clean negatives.  $\mathrm{SU}(3)$ has no natural
matter-coupled lattice gauge theory---the Clebsch--Gordan intertwiner
obstruction caps the natural gauge reach (a face of the Coulomb/HO
asymmetry).  The scalar graph photon misses seven of the eight selection
rules because it carries no angular quantum numbers; recovering them is a
calibration cost ($1/(4\pi)$ per loop), not graph-native content.  Twelve
single-mechanism derivations of the $\alpha$ combination rule were
eliminated.  The gravity sector contributes five recorded negatives---the
M\"obius finite-spacing artifact, the failed IR-over-count cure, the
analytically-supplied (not extracted) tip coefficient, the Wald
bookkeeping factor of~2, and the non-area-law wedge entropy---each kept
on the record so it is not re-derived.

\paragraph{Positioning.}
Read as a whole, the Group~5 arc is a structural map of how far the
spectral geometry of $S^3$ reaches into electromagnetism, the gauge
sector, and gravity, with the boundary between forced structure and free
coupling drawn explicitly at every step.  The gauge group structure, the
selection-rule partition, the QED period tower, and the two-term gravity
action are the structural assets;\ the fine-structure constant, the QED
multi-loop constants, and the gravity cutoff are the calibration side.
The discrete and continuous descriptions are dual, with no ontological
priority asserted for either.  The arc is not a derivation of $\alpha$,
not a derivation of the Standard-Model gauge group, and not a theory of
quantum gravity;\ it is a disciplined account of which structures the
framework forces and which numbers it must be told.

%% ========================================================================
%% BIBLIOGRAPHY
%% ========================================================================
\begin{thebibliography}{99}

\bibitem{loutey_paper0}
J.~Loutey,
``Angular Momentum as Information Geometry:\ Isotropic Packing and the
Universal Angular Cross-Section,''
GeoVac Paper~0 (2026).

\bibitem{loutey_paper2}
J.~Loutey,
``The Fine Structure Constant as a Spectral Coincidence on the Hopf
Fibration,''
GeoVac Paper~2 (2026).
[The combination rule $K=\pi(B+F-\Delta)$ is recorded strictly as an
Observation; see \S\ref{sec:alpha}.]

\bibitem{loutey_paper7}
J.~Loutey,
``The Dimensionless Vacuum:\ Recovering the Schr\"odinger Equation from
Scale-Invariant Graph Topology,''
GeoVac Paper~7 (2026).

\bibitem{loutey_paper18}
J.~Loutey,
``Spectral-Geometric Exchange Constants:\ Projecting Discrete Spectra
onto Continuous Manifolds,''
GeoVac Paper~18 (2026).

\bibitem{loutey_paper24}
J.~Loutey,
``The Bargmann--Segal Lattice:\ $\pi$-Free Discretization of the
Harmonic Oscillator on $S^5$,''
GeoVac Paper~24 (2026).

\bibitem{loutey_paper25}
J.~Loutey,
``The Hopf Graph as a Lattice Gauge Structure:\ A Framework Observation,''
GeoVac Paper~25 (2026).

\bibitem{loutey_paper28}
J.~Loutey,
``QED on $S^3$:\ Spectral Geometry of Perturbative Transcendentals,''
GeoVac Paper~28 (2026).

\bibitem{loutey_paper30}
J.~Loutey,
``$\mathrm{SU}(2)$ Wilson Lattice Gauge Structure on the GeoVac Hopf
Graph:\ The Non-Abelian Sibling of Paper~25,''
GeoVac Paper~30 (2026).

\bibitem{loutey_paper32}
J.~Loutey,
``The GeoVac Spectral Triple,''
GeoVac Paper~32 (2026).

\bibitem{loutey_paper33}
J.~Loutey,
``The $1{+}6{+}1$ Partition of QED Selection Rules on $S^3$:\ Graph
Topology, Angular Momentum, and the Dirac Spinor Phase Constraint,''
GeoVac Paper~33 (2026).

\bibitem{loutey_paper34}
J.~Loutey,
``GeoVac as a Two-Layer Framework:\ Projection Taxonomy and the
Empirical Matches Catalog,''
GeoVac Paper~34 (2026).

\bibitem{loutey_paper36}
J.~Loutey,
``Bound-State QED on $S^3$:\ One-Loop Closure of the Hydrogen Lamb
Shift,''
GeoVac Paper~36 (2026).

\bibitem{loutey_paper38}
J.~Loutey,
``State-space Gromov--Hausdorff convergence of truncated
Camporesi--Higuchi spectral triples on $S^3$,''
GeoVac Paper~38 (2026).

\bibitem{loutey_paper41}
J.~Loutey,
``Wilson $U(1)$ Lattice Gauge on the Dirac-Rule-B Graph:\ Structural
Compatibility with 3D Compact $U(1)$ on a Non-Cubic Substrate, Including
the Polyakov Monopole-Density Signature,''
GeoVac Paper~41 (2026).

\bibitem{loutey_paper51}
J.~Loutey,
``Gravity from the GeoVac Spectral Action:\ Continuum Results and the
Discrete-Substrate Program,''
GeoVac Paper~51 (2026).

\bibitem{marcolli_vs}
M.~Marcolli and W.~D.~van Suijlekom,
``Gauge networks in noncommutative geometry,''
J.~Geom.\ Phys.\ \textbf{75}, 71 (2014); arXiv:1301.3480.

\bibitem{chamseddine_connes}
A.~H.~Chamseddine and A.~Connes,
``The Spectral Action Principle,''
Commun.\ Math.\ Phys.\ \textbf{186}, 731 (1997).

\end{thebibliography}

\end{document}
